% =============================================================================
%  inertia_symmetry_obstruction.tex
%
%  A post-mortem on the *inertia-group* route to \Cref{thm:main}
%  (unramifiedness of \cP^{[d]_G} at \eta_C).  This was the first approach we
%  tried: translate the geometric hypotheses (H1), (H2) into constraints on a
%  single inertia group J sitting inside the big Galois group \Gamma = G^d
%  \rtimes S_d, and try to conclude by Galois combinatorics -- it does not
%  work.
%
%  The note holds TWO self-contained summaries of the same dead end, at two
%  altitudes:
%    Part I  (short form)  -- the abstract Goursat picture: (H1)+(H2) fix two
%                             projections of J, a diagonal subgroup escapes.
%    Part II (long form)   -- the explicit local-model computation from the
%                             working session of 2026-05-14: J is computed
%                             outright at d.p_infty, the diagonal is shown to
%                             be populated, an apparent counterexample appears,
%                             and the precise role of higher ramification is
%                             pinned down.
%
%  Compile with:  latexmk -pdf inertia_symmetry_obstruction.tex   (uses biber)
% =============================================================================
\documentclass[11pt]{article}
\usepackage{geometry}\geometry{margin=1in}
\usepackage[T1]{fontenc}
\usepackage{ifthen}
\usepackage{booktabs}

\setlength{\parindent}{1em}
\setlength{\parskip}{0.5em}

\input{includes}
\input{MacrosAndFigures}

% --- local fallbacks for symbols not in the thesis macro file ---------------
\providecommand{\PP}{\mathbb{P}}
\providecommand{\fm}{\mathfrak{m}}

\usepackage{csquotes}
\usepackage[
backend=biber,
style=alphabetic,
sorting=ynt
]{biblatex}
\addbibresource{Bib.bib}

\newboolean{ThesisIntro}
\setboolean{ThesisIntro}{false}

\title{The inertia group is under-constrained\\[2pt]
       {\large why the pure-symmetry route to ``$\cP^{[d]_G}$ unramified at
        $\eta_C$'' stalls, and where higher ramification has to enter}}
\author{Assaf Marzan}
\date{\today}

\begin{document}
\maketitle

\noindent\textbf{Purpose.}
This note records, and explains the failure of, the first approach we tried
for the unramifiedness theorem
\[
\textbf{(H1)} + \textbf{(H2)} + \textbf{(H3)}\ \Longrightarrow\
\cP^{[d]_G}\to U^{(d)}\ \text{unramified at}\ \eta_C .
\]
The approach is attractive because it is purely formal: it stays inside the
Galois group of the master diagram, encodes the whole theorem as a single
containment of inertia groups, and never invokes local ramification theory.
The outcome is sharp and negative.

The note contains \emph{two} summaries of the same dead end, which may be read
independently:
\begin{description}
  \item[\Cref{part:short} (short form).] The abstract reason it fails: (H1) and
        (H2) determine only two projections of the relevant inertia group, and
        by Goursat's lemma a \emph{diagonal} subgroup is left unconstrained --
        precisely the one that breaks the conclusion.
  \item[\Cref{part:long} (long form).] The detailed working session of
        2026-05-14: we computed the inertia group \emph{outright} in the local
        model at the degenerate point $d\cdot p_\infty$, found the diagonal is
        genuinely populated, hit an apparent counterexample to the theorem as
        first stated, and isolated exactly what higher ramification must (and
        must not) be asked to do.
\end{description}
Throughout we assume the setup, notation and master diagram of the standalone
proof note, recalling only what is needed. The common notation is fixed once,
in \Cref{sec:recall}, and used by both parts.

% #############################################################################
\section{Common set-up}\label{sec:recall}
% #############################################################################

Recall the four symmetric quotients
$\cP^{[d]_G},\ \cP^{[d]_H},\ (\cP_{G/H})^{(d)},\ (\cP_{G/H})^{[d]}$ and the master
diagram with Cartesian upper square. All of them are quotients of the single
connected $G^d$-cover $\cP^{\times d}\to U^d$, so the whole picture is governed
by one Galois extension. Writing
\[
K := K(U^{(d)}) = K(C^{(d)}),\qquad N := K(\cP^{\times d}),
\]
the extension $N/K$ is Galois with group
\[
\Gamma \;:=\; G^d \rtimes S_d ,
\]
the $S_d$ permuting the $d$ factors of the diagonal $G^d$. With
$\KG=\ker(G^d\xrightarrow{\,\Sigma\,}G)$ and $K_H=\ker(H^d\to H)$, the five fields
are the fixed fields of subgroups of $\Gamma$:
\[
\begin{array}{llll}
L  := K(\cP^{[d]_G})        & \leftrightarrow & \Xi := \KG\rtimes S_d,
   & \Gal(L/K)\cong G,\\[2pt]
M  := K\big((\cP_{G/H})^{[d]}\big) & \leftrightarrow & (H^d+\KG)\rtimes S_d,
   & \Gal(M/K)\cong G/H,\\[2pt]
K' := K\big((\cP_{G/H})^{(d)}\big) = K(C'^{(d)}) & \leftrightarrow & H^d\rtimes S_d,
   & \\[2pt]
L_H := K(\cP^{[d]_H})       & \leftrightarrow & K_H\rtimes S_d. &
\end{array}
\]
The diagonal map $L/K$ has group $G$ via summation $\Sigma:\Gamma\to G$,
$((g_i),\sigma)\mapsto\sum_i g_i$; the conclusion we want is that a certain
inertia group lands in $\ker\Sigma=\KG\rtimes S_d$. The four fields $M,K',L,L_H$
sit in a \emph{parallelogram} over $M$: one checks $M\subseteq K'$,
$L\cap K' = M$ and $L\cdot K' = L_H$, so $L/M$ and $K'/M$ are linearly disjoint
with compositum $L_H$ and
\begin{equation}\label{eq:parallelogram}
\Gal(L_H/M)\;=\;\underbrace{\Gal(L/M)}_{A}\ \oplus\ \underbrace{\Gal(K'/M)}_{B}.
\end{equation}

Fix a discrete valuation $\widetilde W$ on $N$ extending $V_C$, chosen so that
$\widetilde W|_{K'} = V_{C'}$ (possible because $\Gamma$ acts transitively on the
valuations above $V_C$), and let
\[
J \;:=\; I\big(\widetilde W / V_C\big)\;\subseteq\;\Gamma
\]
be its inertia group. The conclusion of the theorem is exactly the single
containment
\begin{equation}\label{eq:goal}
J \;\subseteq\; \Xi=\KG\rtimes S_d ,
\qquad\text{equivalently}\qquad
\textstyle\sum_i g_i = 0\ \text{in}\ G\ \ \forall\,\big((g_i),\sigma\big)\in J ,
\end{equation}
since the image of $J$ in $G$ under $\Sigma$ is the inertia of $\widetilde W|_L$
over $V_C$.

% #############################################################################
% #############################################################################
\part{Short form: the Goursat gap}\label{part:short}
% #############################################################################
% #############################################################################

\section{Three projections of one inertia group}\label{sec:attempt}

We are forced up to $N$ by the \emph{hypotheses}, not the conclusion: (H1)
lives in $\Gal(M/K)=G/H$, (H2) in $\Gal(L_H/K')=H$, the goal in
$\Gal(L/K)=G$ -- three different quotients of $\Gamma$ with no common Galois
field below $N$. So we read three projections off the single inertia group $J$:
\begin{center}
\begin{tabular}{lll}
\toprule
projection of $J$ & equals inertia of & status\\
\midrule
$J\xrightarrow{\ \Sigma\ }G$ & $\widetilde W|_L$ over $V_C$ & the goal\\
$J\xrightarrow{\ \Sigma\bmod H\ }G/H$ & $\widetilde W|_M$ over $V_C$ & trivial by (H1)\\
$J\cap(H^d\rtimes S_d)\xrightarrow{\ \Sigma\ }H$ & $\widetilde W|_{L_H}$ over $V_{C'}$ & trivial by (H2)\\
\bottomrule
\end{tabular}
\end{center}
In symbols,
\begin{align}
\text{(H1)}:&\quad \textstyle\sum_i g_i \in H\quad\forall\,\big((g_i),\sigma\big)\in J ;
\label{eq:H1}\\
\text{(H2)}:&\quad J\cap\big(H^d\rtimes S_d\big)\subseteq K_H\rtimes S_d,\ \ \text{i.e.}\
\textstyle\sum_i h_i=0\ \text{whenever}\ \big((h_i),\sigma\big)\in J,\ h_i\in H.
\label{eq:H2}
\end{align}
By \cref{eq:H1} every element of $J$ has $(g_i)\in H^d+\KG$; writing
$g_i=h_i+k_i$ with $(h_i)\in H^d$, $(k_i)\in\KG$ (so $\sum_i k_i=0$), the goal
\cref{eq:goal} becomes
\begin{equation}\label{eq:goal2}
\textstyle\sum_i h_i = 0\ \text{in}\ H\qquad\forall\,\big((g_i),\sigma\big)\in J .
\end{equation}

\section{Where the symmetry runs out}\label{sec:stall}

Hypothesis \cref{eq:H2} delivers \cref{eq:goal2} \emph{only} for elements that
already lie in $H^d$ (the case $(k_i)=0$). For a general element with
$(k_i)\neq0$ neither hypothesis says anything: \cref{eq:H2} is about elements
\emph{inside} $H^d\rtimes S_d$, and \cref{eq:H1} only sees the $G/H$-component,
which is vacuous here.

Concretely, the obstruction is a \textbf{diagonal inertia subgroup}. Nothing in
\cref{eq:H1,eq:H2} prevents
\begin{equation}\label{eq:diagonal}
J \;\supseteq\; \big\{\, \big(h,\ \kappa(h)\big) \;:\; h\in H_0 \,\big\}
\;\subseteq\; H \oplus K_{G/H},
\end{equation}
glued along a homomorphism $\kappa:H_0\to K_{G/H}$: it projects trivially to
$G/H$ (so \cref{eq:H1} holds) and meets $H^d\rtimes S_d$ only in $\kappa(h)=0$
(so \cref{eq:H2} holds vacuously), yet has $\sum_i h_i\neq0$, breaking
\cref{eq:goal2}.

\paragraph{The structural reason.}
The symmetry available to us -- the two quotient/restriction maps coming from
$H\trianglelefteq G$ -- only resolves $J$ along the two coordinate axes of the
split sequence $0\to K_H\to\KG\to K_{G/H}\to0$. An inertia group need not
respect that splitting; it can sit diagonally across it. Knowing two coordinate
projections of a subgroup of a product never determines the subgroup -- this is
exactly the content of Goursat's lemma -- so there is simply \emph{not enough
symmetry} in the discrete inertia data alone to collapse the diagonal. The same
point inside the parallelogram \cref{eq:parallelogram}: an inertia subgroup
$\widetilde I\subseteq A\oplus B$ compatible with (H1),(H2) can be either the
harmless $B$ (then $L/M$ unramified, theorem holds) or the fatal diagonal
$\Delta$ (then $L/M$ ramified, theorem fails), and the hypotheses cannot tell
them apart.

% #############################################################################
% #############################################################################
\part{Long form: the explicit local computation}\label{part:long}
% #############################################################################
% #############################################################################

\noindent The short form leaves open whether the diagonal of \cref{eq:diagonal}
is a real possibility or an artefact of weak bookkeeping. In the session of
2026-05-14 we settled this by computing $J$ outright in the local model. The
answer: the diagonal is genuinely there, the theorem as first stated has a
counterexample, and higher ramification enters in a specific, limited way.

\section{Computing \texorpdfstring{$J$}{J} in the local model}\label{sec:localmodel}

Assume $\cP\to U$ is totally ramified at $p_\infty$ with inertia
$G_0\subseteq G$ (else replace $G$ by $G_0$). Choose a uniformiser $t$ on $C$
and $u_j$ on $\overline\cP$ with $t_j=u_j^{|G_0|}+\cdots$. The symmetric local
ring is $k[[e_1,\dots,e_d]]$ in the elementary symmetric functions of the
$t_j$, and the geometrically correct valuation $V_C$ is the \emph{weighted}
blowup with
\[
w(e_i)=i ,\qquad \widetilde w(u_j)=1 .
\]
(With the ordinary unweighted blowup $V_C$ does not even extend to $V_{C'}$.)
Decompose $J$ along $K\subseteq K(C^d)\subseteq N$.

\paragraph{The $S_d$-part is trivial.}
The residue field of $V_C$ is that of the weighted projective space
$\PP(1,2,\dots,d)$, and $S_d$ acts faithfully on the residue field of the lift
to $K(C^d)$ (it permutes the $\PP^{d-1}$-coordinates). So no nontrivial
permutation lies in inertia: $J\to S_d$ is trivial, i.e.\ $J\subseteq G^d$.

\paragraph{The $G^d$-part, tame case: $J=G_{\mathrm{diag}}$.}
In the tame case $g_j$ acts by $u_j\mapsto\zeta^{g_j}u_j$ ($\zeta$ a primitive
$|G_0|$-th root of unity). On the residue field $k(\PP^{d-1}_u)$ the tuple
$(g_1,\dots,g_d)$ acts by
$(u_1:\dots:u_d)\mapsto(\zeta^{g_1}u_1:\dots:\zeta^{g_d}u_d)$, trivial in
$\mathrm{PGL}_d$ iff $g_1=\dots=g_d$. Hence
\begin{equation}\label{eq:Jdiag}
\boxed{\,J \;=\; G_{\mathrm{diag}}\;=\;\{((g,g,\dots,g),e):g\in G_0\}\,}\qquad(\text{tame, totally ramified}).
\end{equation}

\paragraph{The $G^d$-part, wild case: $J\supseteq G^d$.}
If $G$ acts wildly with conductor $m\ge1$ then $g_j(u_j)=u_j+a\,u_j^{m+1}+\cdots$,
and the correction has $\widetilde w$-value $\ge m+1>0$, hence \emph{vanishes in
the residue field}. So \emph{every} tuple in $G^d$ acts trivially on the residue
of $\widetilde W$: $J\supseteq G^d$, strictly larger than the diagonal.

\paragraph{Reading off the image in $G$.}
Under summation, $G_{\mathrm{diag}}\to G$ is $(g,\dots,g)\mapsto d\cdot g$, so
\begin{equation}\label{eq:image}
\operatorname{Im}(J\to G)\;=\;d\cdot G_0 ,
\end{equation}
trivial \emph{iff $|G_0|\mid d$}. The leading-order inertia at $\eta_C$ is thus
the \emph{cardinality} condition $|G_0|\mid d$ -- nothing about the depth of
ramification.

\section{The diagonal is genuinely there: \texorpdfstring{$K'/M$}{K'/M} is ramified}\label{sec:Kprime}

One might hope the diagonal is excluded by some automatic fact. It is not. The
would-be saving hypothesis is
\[
\textnormal{(H3$'$)}\qquad K'/M\ \text{unramified at}\ V_{C'},
\quad\text{i.e.}\quad J\cap\Gal(N/M)\subseteq H^d\rtimes S_d,
\]
which would force $(k_i)=0$ in \cref{eq:goal2} and finish the proof. But the
same weighted-blowup computation gives the inertia of $K'/M$ explicitly. With
$E=G/H$ acting on $C'$ by $u\mapsto\zeta u$, an element
$\vec a\in E^d_+=\ker(\Sigma:E^d\to E)$ acts by
$\epsilon_i\mapsto e_i(\zeta^{a_1}u_1,\dots)$; the \emph{diagonal}
$\vec a=(a,\dots,a)$ scales $\epsilon_i\mapsto\zeta^{ia}\epsilon_i$ and so fixes
every weight-zero coordinate $\xi_i=\epsilon_i/\epsilon_1^i$, while non-diagonal
$\vec a$ act nontrivially. Hence
\begin{equation}\label{eq:gcd}
I_{V_{C'}}(K'/M)\;=\;E_{\mathrm{diag}}\cap E^d_+\;\cong\;\ker(d:E\to E)\;\cong\;\Z/\gcd(d,|E|),
\end{equation}
nontrivial as soon as $\gcd(d,|E|)>1$ -- e.g.\ \emph{totally ramified} when
$d=|E|=2$. Two features make this fatal:
\begin{itemize}
  \item \textbf{It depends only on $E$ and $d$, not on $G$.} The map
        $q:(\cP_{G/H})^{(d)}\to(\cP_{G/H})^{[d]}$ sees only $\cP_{G/H}$, so
        \emph{no} hypothesis on the ambient $G$ -- cyclic, prime-power, or a
        ramification bound on $\cP$ at $p_\infty$ -- can change \cref{eq:gcd}.
  \item \textbf{The available lemma does not reach $V_{C'}$.} $q$ is known
        unramified only at codimension-$1$ boundary valuations of $C'^{(d)}$,
        whereas $V_{C'}$ sits over the codimension-$d$ point $d\cdot p'_\infty$,
        where all $d$ points collide at the branch point and the diagonal
        $E$-action stabilises the configuration most strongly.
\end{itemize}

\paragraph{An apparent counterexample to the theorem-as-stated.}
Pushing through to $L/K$ produced not just a stalled proof but a concrete
obstruction. Take $\operatorname{char}\neq2$, $G=\Z/4$, $H=\Z/2$, $E=\Z/2$,
$d=2$, and $\cP\to U$ tamely totally ramified of index $4$ at $p_\infty$. Then
(H3) holds ($G$ cyclic $2^2$), the ramification bound holds vacuously (tame:
$G_1=0$, so $G_d=0$ for all $d$), and one checks (H1), (H2) both hold. Yet
$L=K(\sqrt[4]{e_2})$ with $v_C(e_2)=2$, so $v_L(\sqrt[4]{e_2})=\tfrac12$ and
$L/K$ is \emph{ramified} at $V_C$ with inertia of order $2$ -- exactly the
predicted $d\cdot G_0=2\cdot\Z/4=H$ of \cref{eq:image}. The missing hypothesis
is visibly $|G_0|\mid d$, which fails here ($4\nmid2$).

\section{Why higher ramification is the right -- and necessary -- language}\label{sec:higher}

The inertia group is the bottom step $G^0=G_0$ of the ramification filtration;
it remembers \emph{whether} a place ramifies but discards the \emph{jumps}. The
diagonal subgroup is invisible to (H1), (H2) precisely because those are
statements about inertia, and inertia cannot separate the $H$-jump from the
$G/H$-jump of a diagonally entangled extension. Two points pin down the role of
higher ramification:

\begin{enumerate}
  \item \textbf{Upper numbering is the natural bookkeeping.} The argument keeps
        passing to quotients $G\to G/H$, and the upper-numbering filtration is
        functorial for quotients via Herbrand, $(G/H)^v=G^v\cdot H/H$, whereas
        lower numbering is functorial for subgroups. So the bound should be read
        in \emph{upper} numbering, $G^d=\{1\}$. This is exactly what makes the
        eventual proof (Herbrand-convexity) work: for cyclic $G=\Z/p^r$ the
        jumps form a single chain and
        \[
        \varphi_G(m_r)=\varphi_{G/H}(m_{r-s})+\frac{\varphi_H(m_r)-m_{r-s}}{p^{r-s}}
        \]
        is a \emph{convex combination}, so (H1),(H2) $\Rightarrow\varphi_G(m_r)<d$.
        Convexity bounds the mixed jump by the two pure jumps \emph{no matter
        how the extension is entangled} -- the very diagonal that defeated the
        inertia argument.
  \item \textbf{But reparameterising alone is not the fix.} The leading-order
        obstruction \cref{eq:image} depends only on the \emph{cardinality}
        $|G_0|=|G^0|$ relative to $d$, and upper vs.\ lower numbering are two
        parameterisations of the \emph{same} filtration -- they do not change
        $G_0$, nor the diagonal $J=G_{\mathrm{diag}}$ that summation sends to
        $d\cdot G_0$. The higher filtration \emph{augments} the leading-order
        picture (it controls the next-order action on tangent and higher
        infinitesimal data at $p'_\infty$, hence the codimension at which the
        ramification of $\cP$ becomes visible in $C^{(d)}$); it does not by
        itself collapse the leading-order diagonal. In the genuinely useful
        regime the bound $G^d=\{1\}$ is non-vacuous and is consumed by the
        Herbrand-convexity estimate, not by a leading-order count.
\end{enumerate}

\section{Summary}\label{sec:summary}

\begin{itemize}
  \item \textbf{The route.} Encode the theorem as $J\subseteq\KG\rtimes S_d$ for
        the inertia group $J\subseteq\Gamma=G^d\rtimes S_d$, and read three
        projections of $J$ off (H1), (H2) and the goal.
  \item \textbf{Why it stalls (short form).} (H1), (H2) fix only two
        projections; by Goursat a diagonal subgroup $\Delta$ is left
        unconstrained, and the hypotheses cannot tell the harmless inertia $B$
        from the fatal $\Delta$.
  \item \textbf{The diagonal is real (long form).} The explicit weighted-blowup
        computation gives $J=G_{\mathrm{diag}}$ (tame) or $J\supseteq G^d$
        (wild), with $\operatorname{Im}(J\to G)=d\cdot G_0$; and $K'/M$ has
        inertia $\Z/\gcd(d,|E|)$, depending only on $E,d$. No hypothesis on $G$,
        and no ramification bound on $\cP$, removes it; the theorem as first
        stated even has a tame counterexample, missing the hypothesis
        $|G_0|\mid d$.
  \item \textbf{The moral.} Discrete inertia data is too coarse. The missing
        rigidity lives in the \emph{higher ramification filtration}: read the
        bound in upper numbering and feed it into the Herbrand-convexity lemma,
        which is how the successful proofs of the unramifiedness theorem close
        the gap. The
        inertia/symmetry argument is non-load-bearing -- it cannot reach the
        conclusion on its own.
\end{itemize}

\printbibliography

\end{document}
